\section{Espacios Vectoriales}

\textbf{Objetivos del Capítulo}

\begin{itemize}
\item
Aprender la definici\'on de un espacio vectorial real.

\item
Estudiar el concepto de subespacio vectorial.

\item
Familiarizarse con la operación de combinación lineal y el concepto de espacio generado.

\item
Conocer la propiedad de independencia lineal.

\item
Profundizar en el conjunto mínimo de vectores con los que se puede generar un espacio vectorial. Para definir el concepto de bases y dimensiones de un espacio vectorial.

\item
Expresar vectores con bases diferentes.

\item
Aprender a definir conceptos relacionados con subespacios vectoriales formados a partir de las filas y las columnas de matrices.

\end{itemize}



\subsection{Definición y Propiedades Básicas}

Durante toda esta secci\'on, denotaremos por $K$ un \emph{cuerpo}, esto es de un conjunto con dos leyes de composici\'on, denotadas $+$ y $\cdot$, y dos elementos especiales 0 y 1, distintos entre s\'i, que cumplen las siguientes
propiedades:

\begin{enumerate}
    \item $+$ y $\cdot$ son asociativas;
    \item $+$ y $\cdot$ son conmutativas;
    \item $\cdot$ es distributiva con respecto a $+$;
    \item $0$ es el neutro de $+$ y $1$ es el neutro de $\cdot$;
    \item para todo $a \in K$, existe $a'\in  K$ tal que $a + a'=0$ (usualmente $a'$ se denota por $-a$);
    \item para todo $a \in K\setminus \{0\}$, existe $a \in K \setminus \{0\}$ tal que $a \cdot a' = 1$ (usualmente $a''$ se denota $a^{-1}$ o $1/a$).
\end{enumerate}
Si bien esta definici\'on es general, durante el curso consideraremos $K$ igual a $\Q, \R$ o $\C$.

Sea $(K, +, \cdot)$ un cuerpo. Comencemos con la definici\'on de un espacio vectorial.
\begin{defin}

Un $K-$espacio vectorial (o espacio vectorial sobre $K$) es un
conjunto no vac\'io $V$ provisto de dos operaciones binarias 
\begin{align*}
        + : V \times V &\to V &\cdot : K \times V &\to V\\
 (x, y) &\mapsto x + y&  (\alpha, y)& \mapsto \alpha \cdot y
\end{align*}
llamadas respectivamente suma y multiplicación por escalar y que satisfacen las siguientes propiedades:

\begin{enumerate}

\item
$+$ es asociativa, es decir $(x+ y)+ z = x+ (y+ z)$ para todo $x, y, z \in V$.

\item
$+$ es conmutativa, es decir $x + y = y + x$ para todo $x, y \in V$.

\item
$+$ tiene un elemento neutro $0_V$ , es decir $x + 0_V = x = 0_V + x$ para todo
$x \in V$.

\item
$+$ tiene inversos, es decir, para todo $x \in V$ , existe $x' \in V$ tal que $x + x' =0_V$.

\item
$\alpha \cdot  (\beta \cdot  x) = (\alpha \cdot \beta) \cdot  x$ para todo $\alpha , \beta \in K$ y todo $x \in V$.

\item
$(\alpha + \beta) \cdot x = (\alpha \cdot x) + (\beta \cdot x)$ para todo $\alpha,\beta \in K$ y todo $x \in V$.

\item
$\alpha \cdot (x + y) = (\alpha \cdot x) + (\alpha \cdot y)$ para todo $\alpha \in K$ y todo $x, y \in V$.

\item
$1 \cdot x = x$ para todo $x \in V$.
\end{enumerate}
\end{defin}
Los elementos de un espacio vectorial se llaman vectores.
Denotaremos por $(V,+,\cdot, K)$ al $K-$espacio vectorial $V$.



\begin{ejem}
\begin{enumerate}
\item
$\mathbb{R}^2$ es un espacio vectorial sobre $\R$.

\item
$\mathbb{R}^n$ es un espacio vectorial sobre $\R$.

\item
$V= \{ \textbf{0} \}$ es un espacio vectorial sobre cualquier cuerpo.

\item
$V= \{ 1 \}$ No es un espacio vectorial.

\item
$V= \{ (x,y)\in \R^2: y = mx \}$ es un espacio vectorial sobre $\R$.


\item
$V = M_{m\times n}(\R)$ el conjunto de matrices $m\times n$ con coeficientes en $\R$ es
un $\R-$espacio vectorial para la suma de matrices y la  multiplicaci\'on por escalares.

\item $\C$ es un $\R-$espacio vectorial con la suma de complejos y la multiplicaci\'on por escalar dada por $\alpha(a+bi) = \alpha a+\alpha bi$ para todo $\alpha \in \R$ y $(a+bi)\in \C$.
$\C$ es un $C-$espacio vectorial con la suma de complejos y la  la multiplicaci\'on de complejos.
Estos espacios vectoriales son distintos.
\item Definimos el conjunto de todos los polinomios con coeficientes en $K$ con indeterminada $x$ como
\[K[x]=\left\{p(x)=\sum_{i=0}^m a_i x^i\quad\colon a_i\in K, \quad m \in \N\right\},\]
con las operaciones adici\'on y ponderaci\'on definidas por
\begin{enumerate}
    \item $\sum_{i=0}^s a_i x^i+\sum_{i=0}^t b_i x^i=\sum_{i=0}^m (a_i+b_i) x^i$, donde $m=\max\{s,t\}$.
    \item $\alpha\left(\sum_{i=0}^m a_i x^i\right)=\sum_{i=0}^m (\alpha a_i) x^i$.
\end{enumerate}
Con estas operaciones $(K[x],+, \cdot, K)$ es un espacio vectorial.
\item Definimos el conjunto 
\[K_n[x]=\left\{p(x)\in K[x]\colon \deg(p(x))\leq n\right\}.\]
Este conjunto con las operaciones anteriores, es un espacio vectorial sobre $K$.
\item Sea $A\subset \R$ no vac\'io. El conjunto de funciones 
\[\mathcal{F}(A,\R):=\{f:A\to \R\colon f \text{ funci\'on}\},\]
con la suma y poderaci\'on usuales(punto a punto), es un $\R-$espacio vectorial.
\item Sea $A\subset \R$ no vac\'io. El conjunto de funciones continuas
\[\mathcal{C}(A,\R):=\{f:A\to \R\colon f \text{ funci\'on continua}\},\]
con la suma y poderaci\'on usuales, es un $\R-$espacio vectorial.
\item Sea $A\subset \R$ no vac\'io. El conjunto de funciones continuas
\[\mathcal{D}(A,\R):=\{f:A\to \R\colon f \text{ funci\'on derivable}\},\]
con la suma y poderaci\'on usuales, es un $\R-$espacio vectorial.
\end{enumerate}
\end{ejem}

\begin{ejer}
Verifique que los ejemplos anteriores son efectivamente espacios vectoriales sobre los cuerpos correspondientes.
\end{ejer}



\begin{thm}
Sea $V$ un espacio vectorial. Entonces:

\begin{enumerate}
\item
$\alpha \textbf{0} = \textbf{0}$ para todo escalar $\alpha$.

\item
$\textbf{0} \cdot x = \textbf{0}$ para todo $x \in V$.

\item
Si $\alpha x = \textbf{0}$, entonces $\alpha = 0$ o $x = \textbf{0}$ (o ambos).

\item
$(-1)x = -x$ para todo $x \in V$.
\end{enumerate}
\end{thm}
\begin{ejer}
Demuestre las propiedades anteriores.
\end{ejer}
\subsection{Subespacios vectoriales}
\begin{defin}
Se dice que $H$ es un subespacio vectorial de $V$ si $H$ es un subconjunto no vacío de $V$, y $H$ es un espacio vectorial, junto con las operaciones de suma entre vectores y multiplicación por escalar definidas para $V$.
\end{defin}

\begin{thm}
Un subconjunto no vacío de $H$ de un espacio vectorial $V$ es un subespacio de $V$ si se cumplen las siguientes condiciones:

\begin{enumerate}
\item
Si $x \in H$ y $y \in H$, entonces $x + y \in H$.

\item
Si $x \in H$, entonces $\alpha x \in H$ para todo escalar $\alpha$.
\end{enumerate}
\end{thm}

\begin{rmk}
Todo subespacio de un espacio vectorial $V$ contiene al vector $\textbf{0}$.
\end{rmk}

La definici\'on de subespacio se puede reducir a una sola propiedad, que es la siguiente: para todo $\alpha \in K$ y para todo $x, y \in H$, $\alpha x+y \in H$. En efecto, esta propiedad se deduce inmediatamente de (1) y (2). Por otra parte, partiendo de esta propiedad
obtenemos (1) tomando $\alpha = 1$. Con $x = y$ y $\alpha = -1$ obtenemos que $0\in H $, por
lo que podemos tomar $y = 0$ y obtenemos (2). En particular, vemos que todo
subespacio vectorial contiene a $0$.




\begin{ejem}
\begin{enumerate}
\item
$H = \{ \textbf{0} \}$.

\item
Todo espacio vectorial es subespacio de sí mismo.


\item
Sea $V = \mathbb{R}^2$. Un subespacio vectorial de $V$ es $H = \{ (x,y): y = mx \}$.


\item
Demuestre que $H_1 = \{(x,y,z): 2x - y - z = 0 \}$ es un subespacio de $\mathbb{R}^3$
\item Demuestre que $H = \{ (x,y,z): x + 2y + 3z = 0 \}$ es un subespacio vectorial de $\mathbb{R}^3$
\end{enumerate}
\end{ejem}



\begin{thm}
Sean $H_1$ y $H_2$ subespacios de un espacio vectorial $V$. Entonces $H_1 \cap H_2$ es un subespacio vectorial de $V$.
\end{thm}

\begin{ejem}
Sea $V = \mathbb{R}^3$ y sean $H_1 = \{(x,y,z): 2x - y - z = 0 \}$ y $H_2 = \{ (x,y,z): x + 2y + 3z = 0 \}$. Entonces $H_1$ y $H_2$ son subespacios de $\mathbb{R}^3$. $H_1 \cap H_2$ es la intersecci\'on de los dos planos.
\end{ejem}


\subsection{Combinación Lineal e Independencia Lineal}

\begin{defin}
Sean $v_1, v_2, \dots, v_n$ vectores en un espacio vectorial $V$. Entonces cualquier vector, $v$, de la forma: 
$$v = a_1 v_1 + a_2v_2 + \cdots +a_n v_n$$ 

donde $a_1, a_2, \dots a_n$ son escalares, se denomina una combinación lineal de $v_1, v_2, \dots , v_n$.
\end{defin}


\begin{ejem}
Consideremos el vector $\begin{pmatrix} -7 \\ 7 \\ 7  \end{pmatrix} \in \mathbb{R}^3$. 
\\\

Como $\begin{pmatrix} -7 \\ 7 \\ 7  \end{pmatrix} = 2\begin{pmatrix} -1 \\ 2 \\ 4 \end{pmatrix} - \begin{pmatrix} 5 \\ -3 \\ 1 \end{pmatrix}$ 
\\\

luego $\begin{pmatrix} -7 \\ 7 \\ 7  \end{pmatrix}$ es una combinaci\'on lineal de $\begin{pmatrix} -1 \\ 2 \\ 4 \end{pmatrix}$ y $\begin{pmatrix} 5 \\ -3 \\ 1 \end{pmatrix}$.
\end{ejem}


\begin{defin}
Sea $A=\{v_1, v_2, \dots, v_n\}$ subconjunto de vectores en un espacio vectorial $V$. Se dice que el conjunto $A$ es
Linealmente Independiente (L.I.) si 
$$ a_1 v_1 + a_2v_2 + \cdots +a_n v_n=0$$ 
implica que $a_i=0$ para todo $i=1, \dots, n.$
En caso contrario, es decir, existe $i\in \{1,\dots, n\}$ tal que $a_i\neq 0$, el conjunto $A$ se dice linealmente dependiente (L.D.).
\end{defin}

\begin{rmk}
\begin{enumerate}
\item 
Un conjunto de vectores $v_1,v_2, \dots , v_n$ en un espacio vectorial $V$ es linealmente dependiente si y sólo si al menos uno de ellos puede expresarse como una combinaci\'on lineal de los otros.

    \item Todo subconjunto de vectores L. I. es también L.I.
    \item Todo conjunto de vectores que contenga un conjunto de vectores L.D. es también
L.D.
\item Todo conjunto de vectores que contenga al vector nulo es L.D.
\item Si $v$ es un vector no nulo de un espacio vectorial $V$, entonces el conjunto $\{v\}$ es
L.I.
\end{enumerate}
\end{rmk}

\begin{ejem}
Los vectores $\{(1,1,1) , (1,2,2) , (2,1,-1)\}$ , son linealmente independientes, justifique su
respuesta:
Consideremos la ecuación 
\[\alpha (1,1,1 ) + \beta ( 1,2,2 ) + \gamma ( 2,1,-1) = ( 0, 0,0 ) ,\]
que determina el sistema lineal homogéneo 
\begin{align*}
    \alpha + \beta + 2\gamma &= 0,\\ \alpha + 2\beta + \gamma &= 0,\\ \alpha + 2\beta - \gamma &= 0,
\end{align*} cuyas soluciones son $ \alpha  = 0 $, $\beta = 0$ y $\gamma = 0$, por tanto, los
vectores son linealmente independientes.
\end{ejem}

\begin{ejer}
Determine si los siguientes conjuntos son linealmente independientes o lin. dependiente.
\begin{enumerate}
    \item $\{(2, -1), (1,0)\}$ en el $\R-$ espacio vectorial $\R^2$.
    \item $\{(1,1,0), (1,1,1), (0,1, -1)\}$ en el $\R-$ espacio vectorial $\R^3$.
    \item $\{ 1, x -1, x^2 + 2x + 1, x^2\}$ en el $\R-$ espacio vectorial $\R[x]$.
    \item $\left\{\begin{pmatrix} -2 & 1 \\ 1 & 0 \end{pmatrix},\begin{pmatrix} 1 & -1 \\ 0 & 2 \end{pmatrix},\begin{pmatrix} -1 & 0 \\ 1 & 2 \end{pmatrix},\begin{pmatrix} 0 & -3 \\ 1 & 2 \end{pmatrix}\right\}$ en el $\R-$ espacio vectorial $M_{2}(\R).$
\end{enumerate}
\end{ejer}

\begin{ejer}
Sean $\{v_1, v_2, v_3\}$ vectores linealmente independientes en un espacio vectorial $V$. Demuestre que los siguientes conjuntos son linealmente independientes.
\begin{enumerate}
    \item $\{ v_1 + v_2- v_3, v_2+ v_3, 2v_1\}$.
\item $\{ v_1 + v_2, v_1 - 2v_2 + v_3, v_1 - v_3\}$.
\end{enumerate}
\end{ejer}

\begin{ejer}
Determine el valor de $k \in \R$ de modo que los siguientes conjuntos sean:
\begin{enumerate}
    \item $\{ (1, k, 0), (-1,2,1), (2, -1, k) \}$ linealmente independientes.
    \item ${ -x_1 + x_2 + 3x_3, 5x_1- 2x_2 + 9x_3, x_1 - kx_2+ 2kx_3 } $ linealmente dependiente.
\end{enumerate}
\end{ejer}

\begin{ejer}
Escriba el vector $v = 11 - 9x + 3x^2 \in \R_2[x]$ como combinación lineal de los vectores $v_1 = 1 - x $, $v_2 = 2x + x^2$ y $v_3 = 5 - x + 3x^2$.
\end{ejer}

\begin{defin}
Se dice que los vectores $v_1, v_2 \dots , v_n$ de un espacio vectorial $V$ generan a $V$ si todo vector en $V$ se puede escribir como una combinación lineal de los mismos. Es decir, para todo $v \in V$ existen escalares $a_1,a_2, \dots a_n$ tales que
$$v = a_1 v_1 + a_2 v_2 + \cdots + a_n v_n$$
\end{defin}

\begin{ejem}
\begin{enumerate}
\item
Los vectores $e_1 = \begin{pmatrix} 1 \\ 0 \end{pmatrix}$ y $e_2 = \begin{pmatrix} 0 \\ 1 \end{pmatrix}$ generan $\mathbb{R}^2$
\\\

\item
$e_1 = \begin{pmatrix} 1 \\ 0 \\ 0\end{pmatrix}$, $e_2 = \begin{pmatrix} 0 \\ 1 \\ 0 \end{pmatrix}$ y $e_3 = \begin{pmatrix} 0 \\ 0 \\ 1 \end{pmatrix}$ generan $\mathbb{R}^3$
\end{enumerate}
\end{ejem}





\begin{ejem}
\begin{enumerate}
\item
$1,x,x^2, \dots, x^n$ generan a $\mathbb{R}_n[x]$
\\\

\item
$\begin{pmatrix} 1 & 0 \\ 0 & 0 \end{pmatrix}$, $\begin{pmatrix} 0 & 1 \\ 0 & 0 \end{pmatrix}$, $\begin{pmatrix} 0 & 0 \\ 1 & 0 \end{pmatrix}$ y $\begin{pmatrix} 0 & 0 \\ 0 & 1 \end{pmatrix}$ generan a $M_{2}$
\\\

\item
No hay un conjunto finito de vectores que generen $\mathbb{R}[x]$
\end{enumerate}
\end{ejem}


\subsection{Espacio generado.}

Sea $V$ un espacio vectorial sobre $K$.
\begin{defin}
Sea $v_1, v_2 \dots , v_k$ $k$ vectores de un espacio vectorial $V$. El espacio generado por $\{ v_1,v_2, \dots v_k \}$ es el conjunto de combinaciones lineales de $v_1,v_2, \dots ,v_k$. Es decir, 
\begin{center}
   $<\{v_1, v_2, \dots , v_k \}> = \{ v\in V\colon v =  a_1 v_1 + a_2 v_2 + \cdots + a_k v_k \}$  
\end{center}
donde $a_1, a_2, \dots a_k$ son escalares. De manera m\'as general, sea $S$ un subconjunto de $V$, el subespacio generado por $S$ se define como
\[<S>=\{v\in V\colon v=\sum_{i=1}^n a_i v_i,\quad v_i\in S, \quad a_i \in K, \quad n\in \N\},\]
que corresponde al conjunto de todas las
combinaciones lineales finitas de vectores de $S$ con coeficientes en $K$,
\end{defin}

\begin{thm}
Si $v_1, v_2, \dots, v_k$ son vectores en un espacio vectorial $V$, entonces $<\{v_1, v_2, \dots , v_k \}>$ es un subespacio vectorial de $V$.
\end{thm}

\begin{rmk}
\begin{enumerate}
\item Sea $S$ un conjunto de vectores de un espacio vectorial $V$. El subespacio generado por S,
se define como la intersección de todos los subespacios de $V$ que contienen a $S$.
\item $<S>$ es el menor(m\'as pequeño) subespacio generado por $S$.
\item $<\emptyset>=\{0\}$
\end{enumerate}
\end{rmk}



\begin{ejem}
Sean $v_1 = (2,-1,4)$ y $v_2 = (4,1,6)$. Encuentre $H=<\{ v_1,v_2 \}>$
\end{ejem}

\begin{thm}
Sean $v_1, v_2, \dots, v_n, v_{n+1}$; $n+1$ vectores que están en un espacio vectorial $V$. Si $v_1, v_2, \dots , v_n$ generan a $V$, entonces $v_1, v_2, \dots, v_n, v_{n+1}$ también generan a $V$.
\end{thm}




\begin{ejer}
Demuestre el teorema anterior.
\end{ejer}


\begin{thm}
Cualquier conjunto de $n$ vectores linealmente independientes en $\mathbb{R}^n$ genera a $\mathbb{R}^n$.
\end{thm}









\begin{thm}
Un conjunto de $n$ vectores en $\mathbb{R}^m$ es siempre linealmente dependiente si $n > m$.
\end{thm}

\begin{ejem}
Los vectores  $\begin{pmatrix} 2 \\ -3 \\ 4 \end{pmatrix}, \begin{pmatrix} 4 \\ 7 \\ -6 \end{pmatrix}, \begin{pmatrix} 18 \\ -11 \\ 4 \end{pmatrix}, \begin{pmatrix} 2 \\ -7 \\ 3 \end{pmatrix}$ son Ld.
\end{ejem}

\begin{cor}
Un conjunto de vectores linealmente independientes en $\mathbb{R}^n$ contiene a lo más $n$ vectores.
\end{cor}





\begin{thm}
Sean $v_1, v_2, \dots v_n$, $n$ vectores en $\mathbb{R}^n$ y sea $A$ una matriz de $n \times n$ cuyas columnas son $v_1, v_2, \dots v_n$. Entonces, $v_1, v_2, \dots, v_n$ son linealmente independientes si y solo si la única solución al sistema homogéneo $Ax = 0$ es la solución trivial $x = 0$.
\end{thm}

\begin{thm}
Sea $A$ una matriz de $n \times n$. Entonces $\det A \not= 0$ si y solo si las columnas de $A$ son linealmente independientes.
\end{thm}





\begin{ejem}
\begin{enumerate}
\item
Determine si los vectores $\begin{pmatrix} 1 \\ -2 \\ 3 \end{pmatrix}, \begin{pmatrix} 2 \\ -2 \\ 0 \end{pmatrix}, \begin{pmatrix} 0 \\ 1 \\ 7 \end{pmatrix}$ son Li o Ld.
\\\

\item
Determine si los vectores $\begin{pmatrix} 1 \\ -3 \\ 0 \end{pmatrix}, \begin{pmatrix} 3 \\ 0 \\ 4 \end{pmatrix}, \begin{pmatrix} 11 \\ -6 \\ 12 \end{pmatrix}$ son Li o Ld.
\end{enumerate}
\end{ejem}



\subsection{Bases y Dimensión}

\begin{defin}
Un conjunto finito de vectores $\{ v_1,v_2, \dots, v_n \}$ es una base para un espacio vectorial $V$ si 

\begin{enumerate}
\item
$\{ v_1,v_2, \dots, v_n \}$ es linealmente independiente.

\item
$\{ v_1,v_2, \dots, v_n \}$ genera a $V$.
\end{enumerate}
\end{defin}

\begin{ejem}
En $\mathbb{R}^2$ los vectores $e_1 = \begin{pmatrix} 1 \\ 0  \end{pmatrix}$ y $e_2 = \begin{pmatrix} 0 \\ 1 \end{pmatrix}$ son una base para $\mathbb{R}^2$.
\end{ejem}

\begin{rmk}
Todo conjunto de $n$ vectores linealmente independientes en $\mathbb{R}^n$ es una base en $\mathbb{R}^n$.
\end{rmk}



\begin{defin}
En $\mathbb{R}^n$ al conjunto de vectores 
$$e_1 = \begin{pmatrix} 1 \\ 0 \\ \vdots \\ 0 \end{pmatrix}, e_2 = \begin{pmatrix} 0 \\ 1 \\ \vdots \\ 0 \end{pmatrix}, \cdots, e_n = \begin{pmatrix} 0 \\ 0 \\ \vdots \\ 1 \end{pmatrix}$$ 

se le denomina base canónica.
\end{defin}

\begin{ejem}
Encuentre una base para el conjunto de vectores que se encuentran en el plano 
\[\pi = \left\{ \begin{pmatrix} x \\ y \\ z \end{pmatrix} : 2x-y+3z=0 \right\}\]
\end{ejem}





\begin{ejem}
¿Cuál puede ser la base canónica de $\mathbb{R}_3[x]$ y de $M_{2 }(\R)$?
\end{ejem}





\begin{thm}
Si $\{ v_1,v_2, \dots, v_n \}$ es una base para $V$ y si $v \in V$, entonces existe un conjunto único de escalares $c_1, c_2, \dots, c_n$ tales que $v = c_1v_1+c_2v_2 + \cdots c_nv_n$.
\end{thm}

\begin{thm}
Si $u_1,u_2, \dots u_m$ y $v_1,v_2, \dots, v_n$ son bases en un espacio vectorial $V$, entonces $m = n$.
\end{thm}

\begin{defin}
Sea $V$ un espacio vectorial sobre $K$, diremos que la dimensión de $V$ es $n$ si y sólo si tiene
una base de $n$ vectores y escribimos $\dim_K(V)=n$ o simplemente $\dim(V)=n$ y el cuerpo se da por conocido.

De otra manera, $V$ se denomina espacio vectorial de dimensión infinita. Si $V = \{ \textbf{0} \}$, entonces se dice que $V$ tiene dimensión cero.
\end{defin}





\begin{ejem}
\begin{enumerate}
\item
$\dim_\R\mathbb{R}^n = n$


\item
$\dim_\R\mathbb{\R}_n[x] = n+1$


\item
$\dim_\R M_{mn}(\R) = mn$


\item
$\dim_\R\mathbb{\R}[x] = \infty$
\end{enumerate}
\end{ejem}
\begin{ejer}
Exhiba bases que confirmen los números precedentes en los casos de dimensión finita.
\end{ejer}



\begin{thm}
Suponga que dim$V=n$. Si $u_1,u_2, \dots, u_m$ es un conjunto de $m$ vectores linealmente independientes en $V$, entonces $m \le n$.
\end{thm}

\begin{thm}
Sea $H$ un subespacio vectorial de dimensión finita $V$. Entonces $H$ tiene dimensión finita y dim$H \le$dim$V$.
\end{thm}

\begin{ejem}
\begin{enumerate}
\item
$\mathbb{R}^3$ tiene dimensi\'on 3. Luego, los subespacios de $\mathbb{R}^3$ tienen dimensi\'on 3, 2, 1 o 0.
\\\

\item
Encuentre una base y la dimensi\'on para el espacio de soluci\'on $S$ del sistema:
\begin{center}
    $x+2y-z=0$

$2x-y+3z=0$
\end{center}
\end{enumerate}
\end{ejem}





\begin{ejem}
Encuentre una base para el espacio de soluciones $S$ del sistema:
\begin{center}
    $2x-y+3z=0$

$4x-2y+6z=0$

$-6x+3y-9z=0$
\end{center}
\end{ejem}

\begin{thm}
Cualquier conjunto de $n$ vectores linealmente independientes en un espacio vectorial $V$ de dimensión $n$ constituyen una base para $V$.
\end{thm}

\begin{thm}
Sea $V$ un $K-$espacio vectorial finitamente generado. Sea
$S \subset V$ un subconjunto de $V$ LI sobre $K$. Entonces existe una base $B$ de $V$ tal que
$S \subset B$. En otras palabras, todo subconjunto LI se puede completar a una base de
$V $. En particular, todo espacio vectorial finitamente generado posee una base.
\end{thm}

\subsection{Suma de subespacios}
Tal como vimos anteriormente, no todo subconjunto de un espacio vectorial es también un espacio vectorial. Ya vimos como, a partir de conjuntos de vectores, podemos generar subespacios.
Esto se puede aplicar a la uni\'on de subespacios (la cual, nuevamente, no es un
subespacio). Sin embargo, existe una forma más sencilla de interpretar el espacio generado por dos o más subespacios. Se trata de lo siguiente:

\begin{defin}
Sea $V$ un espacio vectorial sobre $K$ y sean $W_1, \dots, W_n \leq V$ subespacios. Definimos la suma de los subespacios $W_i$ y anotamos $\sum^n_
{i=1} W_i$ (o sencillamente $W_1 + W_2$ si $n = 2$) el conjunto
\[\{v_1+v_2+\cdots+v_n \colon v_i\in W_i \text{ para todo } 1\leq i\leq n\}=<\cup_{i=1}^nW_i>.\]
\end{defin}
Veamos ahora que ocurre con la dimensi\'on de un espacio suma.
\begin{thm}
Sean $W_1, W_2$ subespacios vectoriales de dimensi\'on finita de un
$K-$espacio vectorial $V $. Entonces $W_1 + W_2$ es de dimensi\'on finita y
\[\dim_K(W_1+ W_2)=\dim_K(W_1)+\dim_K(W_2)-\dim_K(W_1\cap W_2).\]
\end{thm}



\begin{defin}
Sea $V$ un espacio vectorial y sean $W_1, W_2$ dos subespacios vectoriales de $V$, se define la
\emph{suma directa} de estos subespacios al subespacio $W_1 \oplus W_2$ si y sólo si $ W_1 + W_2= W$ y además $W_1 \cap W_2 = \{ 0\}$.
\end{defin}
\begin{ejem}
Sean los subespacios vectoriales $W_1= \{ (x, y, z) \in \R^3 \colon z = 0\}$ , $W_2 =\{ (x, y, z) \in \R^3 \colon  x = y = 0 \}$
Se puede observar que $W_1$ representa un plano del espacio y $W_2$ una recta del espacio no coincidente con el plano, pudiéndose comprobar que $W_1 \cap W_2 = \{0\}$, y además que $W_1 + W_2 = \R^3$. Por tanto se puede afirmar que $W_1 \oplus W_2 = \R^3$.
\end{ejem}
\section{Espacios vectoriales con producto interno}

\textbf{Objetivos del Capítulo}

\begin{itemize}
\item
Familiarizarse con la operación de producto interno. 


\item
Estudiar los conceptos de ortogonalidad y proyecciones con especto a espacios vectoriales.

\item
Calcular bases ortonormales y proyecciones ortogonales.

\end{itemize}





\subsection{Espacio Euclideos}

\begin{defin}
Un $K$ espacio vectorial $V$ será llamado euclideo si hay una regla que asigna a cada par de
vectores $u, v \in V$ un número complejo, llamado producto escalar designado por $\langle u , v\rangle$ 
\begin{align*}
    \langle \; , \rangle : V
\times V &\to K\\
    (u,v)&\mapsto \langle u , v\rangle
\end{align*}
de manera que se cumplen las siguientes propiedades:
\begin{enumerate}[(a)]
\item Simetría: $\langle u , v\rangle=\overline{\langle v , u\rangle}$ para todo $u, v \in V$ y donde la barra indica conjugaci\'on compleja.
\item Distributividad: $\langle u , v+w\rangle=\langle u , v\rangle+\langle u , w\rangle$ para todo $u, v, w \in V$.
\item Ponderación: $\langle ku , v\rangle=k\langle u , v\rangle$ para todo $u, v \in V$, y todo $k \in K$.
\item Positividad:$\langle u , u\rangle\geq 0$ , para todo $u \in V$.
\end{enumerate}
De (b) y (a) se deduce que:
\begin{enumerate}
    \item $\langle u+v , w\rangle=\langle u , w\rangle+\langle v , w\rangle$ para todo $u, v, w \in V$.
    \item $\langle u , kv\rangle=k\langle u , v\rangle$ para todo $u, v \in V$, y todo $k \in K$.
\end{enumerate}
Utilizando (1) se tiene que $\langle v , w\rangle=\langle 0+v , w\rangle=\langle 0 , w\rangle+\langle v , w\rangle$ con lo cual $\langle 0 , w\rangle=0$ y por la propiedad simétrica también $\langle w, 0\rangle=0$. En particular $\langle v , v\rangle=0$ si $u= 0$.
Por tanto un espacio vectorial real $V$ será llamado euclideo o espacio con \emph{producto interno}
si hay una función $\langle , \rangle: V \times V \to K$ tal que $(u, v) \to \langle u , v\rangle$, y que satisface
todas las propiedades descritas anteriormente.

\end{defin}

\begin{ejem}
\begin{enumerate}
    \item En $\mathcal{F}(\R,\R)$, espacio vectorial de las funciones reales consideremos el subespacio vectorial $\mathcal{C}(I,\R)$
de todas las funciones continuas en $I = [0,1]$. Se define como producto interno habitual 


\begin{align*}
    \langle \; , \rangle: \mathcal{C}(I,\R)\times \mathcal{C}(I,\R) &\to \R\\
    (f, g) &\mapsto \langle f, g\rangle = \int_0^1 f(x) \cdot g(x)\;dx
\end{align*}
    \item En $M_n(\R)$, se define un producto interno habitual: $\langle A, B\rangle = tr(A^tB)$ o bien $\langle A, B\rangle = tr(BA^t)$.
    \item En $\R^n$, el  producto interno más habitual es el llamado producto interno euclidiano o producto punto entre los vectores $x = (x_1, x_2, \dots,x_n), y = (y_1, y_2,\dots,y_n) \in \R^n$ se define como:
\[\langle x, y \rangle=\sum_{i=1}^n x_i y_i.\]
Verificar para $\R^3$.

\item En el espacio vectorial $\R_2[x]$, se define para cada $p(x), q(x) \in \R_2[x]$ el producto interno entre polinomios definido por $\langle p(x), q(x)\rangle = p(0)  q(0) + p(1)  q(1) + p(2)  q(2)$
\end{enumerate}
\end{ejem}






\subsection{Norma, distancia, vector unitario y ángulo entre vectores}



\begin{defin}
Sea $V$ un $K-$espacio vectorial y sea $v \in V$ vector, se define la longitud o norma de $v$, denotada por $\|v\|$, está dada por \[\|v\| = \sqrt{\langle v, v\rangle}.\] 
\end{defin}



\begin{ejem}
\begin{enumerate}
    \item En $\mathbb{R}^3$ la norma de un vector $v = (x,y,z)$ está dada por $\|v\| = \sqrt{x^2+y^2+z^2}$.
    \item La norma del vector $v = (2,-1,3,4,-6) \in \mathbb{R}^5$ es $\sqrt{2^2+(-1)^2+3^2+4^2+(-6)^2} = \sqrt{66}$.
\end{enumerate}
\end{ejem}



\begin{thm}
Sea $V$ un $K-$espacio vectorial con producto interno, entonces
\begin{enumerate}[(a)]
    \item $\|\alpha v\|= |\alpha| \|v\|$ para todo $\alpha\in K$, todo $v\in V$.
    \item $\|v\|>0$, para todo $v\neq 0$ y $\|0_V\|=0$.
    \item $|\langle v, w \rangle|\leq \|v\|\|w\|$ para todo $v,w\in V$ (Desigualdad de Schwarz).
    \item $\|v+w\|\leq \|v\|+\|w\|$ para todo $v,w\in V$(Desigualdad triangular).
\end{enumerate}
\end{thm}

\begin{defin}
Sea $V$ espacio vectorial con un producto interno. Diremos que $u,v\in V$ son \emph{ortogonales} si $\langle v, w \rangle=0$. En cuyo caso escribiremos $v\perp w$.
\end{defin}
\begin{ejem}
En $\R^2$, los vectores $(2,3)$ y $(6,-4)$ son ortogonales.
\end{ejem}


\begin{prop}
\begin{enumerate}
    \item $0_V\perp v$ para todo $v\in V$.
    \item Si $<u,v>=0$ para todo $v\in V$, entonces $u=0_V$.
    \item $u\perp v$, entonces $v\perp u$ para todo $u,v\in V$.
    \item $u\perp v$ y $v\perp w$ \textbf{no implica} que $u\perp w$ para todo $u,v,w\in V$.
\end{enumerate}
\end{prop}

Si $V$ es un $K$ espacio vectorial y $S = \{v_1, \dots, v_n\}$ es una $K$ base de $V$, entonces para cada $v \in V$, existen escalares únicos $a_1,\dots, a_n $ en $K $ tal que todo elemento $v \in V$ se puede escribir como una combinación lineal de los elementos de la base $S$, es decir:
\[ v=a_1v_1+\dots+ a_nv_n\]
Supongamos que en un $V$ espacio vectorial con producto interno $\langle , \rangle$, existe una base $S = \{v_1, \dots, v_n\}$ tal que $\langle v_i, v_j \rangle= 0$ si $i \neq j$, entonces sigue que:
\[a_i=\frac{\langle v, v_i \rangle}{\langle v_i, v_i \rangle}.\]

\begin{defin}
Sea $V$ un $K$ espacio vectorial con producto interno $\langle ,  \rangle$, y consideremos $S = \{v_1, \dots, v_n\}$ tal que $S\subset V$, entonces $S$ se llamará una \emph{base ortogonal} si:
\begin{enumerate}[(a)]
    \item $S$ es una base de $V$.
    \item si $i \neq j$, entonces $\langle v_i, v_j \rangle=0$.
    \item La coordenada respecto de la base ortogonal $a_i=\frac{\langle v, v_i \rangle}{\langle v_i, v_i \rangle}$ se llamará el i-ésimo coeficiente de Fourier del vector $v$.
\end{enumerate}

\end{defin}
\begin{ejer}

Demuestre que el conjunto $\{( 1 , 0 , 0 , ( 0 , 1 , 0), ( 0 , 0, 1)\}$ es una base ortogonal para $\mathbb{R}^3$.

\end{ejer}


\begin{defin}
Un conjunto de vectores es ortonormal si cualquier par de ellos es ortogonal y cada uno tiene longitud $1$.
\end{defin}

Si se define un vector normal o unitario como aquel cuya norma es uno, entonces normalizar
a un vector $v\in V$, consistirá en ponderarlo por el recíproco de su norma construyendo un
nuevo vector que llamaremos el normal asociado a $v$, es decir
\[\hat{v}=\frac{v}{\|v\|}.\]
\begin{ejer}
\begin{enumerate}
    \item Calcule la norma del vector $\left(\frac{1}{2},\frac{1}{2},\frac{1}{2},\frac{1}{2}\right)$.

    \item Normalice la matriz $A=\begin{pmatrix} 1 &-4 \\ 1&1 \end{pmatrix}$ con el producto interno $tr(A^tB)$
\end{enumerate}
\end{ejer}

\begin{defin} En resumen
\begin{enumerate}
    \item $S=\{v_i\}$ es un conjunto normal si $v_i$ tiene norma 1 para todo $i$. Notaci\'on: $\hat{S}$.
    \item $S=\{v_i\}$ es un conjunto ortogonal si $<v_i,v_j>=0$ para todo $i\neq j$. Notaci\'on: ${S}_\perp$.
    \item $S=\{v_i\}$ es un conjunto ortonormal si $<v_i,v_j>=0$ para todo $i\neq j$ y  $<v_i,v_j>=1$ para todo $i= j$. Notaci\'on: $\hat{S}_\perp$.
\end{enumerate}
Si el conjunto es una base, entonces serán bases normal, ortogonal y ortonormal respectivamente.
\end{defin}

%%El conjunto de todos los vectores ortogonales a un conjunto $S$ no vacío, es un s.e.$v$ de $V$ y se le llama subespacio vectorial ortogonal a $S$. Notación: $S^\perp $.
%$S^\perp$ es s.e.v. ortogonal con $<S>$, se dice que $ S^\perp$ y $<S>$ son s.e.v. complementos.

\begin{defin}
\'Angulo entre vectores: Si el ángulo formado por $v$ y $w$ mide $\alpha$, entonces $ cos \alpha =\frac{<v, w>}{\|v\|\| w\|}.$
La distancia entre dos vectores se define como $d(u,v)=\|u-v\|.$
\end{defin}


\begin{thm}
Si $S = \{ v_1,v_2, \dots, v_k  \}$ es un conjunto ortogonal de vectores diferentes de cero, entonces $S$ es linealmente independiente.
\end{thm}















%\begin{itemize}
%\item
%\textbf{Def: }Una matriz $Q$ de $n \times n$ se llama ortogonal si $Q$ es invertible y $Q^{-1} = Q^T$
%\\\

%\item
%\textbf{Teo: }La matriz $Q$ de $n \times n$ es ortogonal si y sólo si las columnas de $Q$ forman una base ortonormal para %$\mathbb{R}^n$.
%\\\

%\item
%\textbf{Ej: }$Q = \begin{pmatrix} \frac{1}{ \sqrt{2} } & -\frac{1}{ \sqrt{6} } & \frac{1}{ \sqrt{3} } \\ \frac{1}{ \sqrt{2} } & \frac{1}{ \sqrt{6} } & -\frac{1}{ \sqrt{3} }   \\ 0 & \frac{2}{ \sqrt{6} } & \frac{1}{ \sqrt{3} }  \end{pmatrix}$
%\end{itemize}
\section[Proyecci\'on ortogonal]{Proyecci\'on ortogonal, distancia a un subespacio y complemento ortogonal.}


\begin{defin}Sea $H$ un subespacio de $\mathbb{R}^n$ con base ortonormal $\{ u_1,u_2, \dots , u_k  \}$. Si $v \in \mathbb{R}^n$, entonces la proyección ortogonal de $v$ sobre $W$, denotada por $P_W(v)$, está dada por:
$P_W( v) = <v , u_1>u_1 + <v , u_2>u_2 + \cdots + <v , u_k>u_k$. En el caso de que la base sea ortogonal, debemos dividir cada sumando por el cuadrado de la norma de $u_i$.
\end{defin}


\begin{ejer}
Encuentre $P_\pi (v)$ donde
 $\pi = \left\{ \begin{pmatrix} x \\ y \\ z \end{pmatrix}: 2x-y+3z=0 \right\}$ y $v = \begin{pmatrix} 3 \\ -2 \\ 4 \end{pmatrix}$
\end{ejer}
\begin{prop}
Sea V un K espacio vectorial con producto interno $< , >$, $H$ subespacio de $V$ y $\{ w_1,w_2, \dots, w_k \}$ una base
ortogonal de W, entonces $P_W(v)=v$ si y s\'olo si $v\in W$.
\end{prop}
\begin{cor}
Sea V un K espacio vectorial con producto interno $< , >$, $W$ subespacio de $V$ y $\{ w_1,w_2, \dots, w_k \}$ una base
ortogonal de W, entonces $P_W\circ P_W=P_W$.
\end{cor}



%\begin{thm}
%Sea $H$ un subespacio de $\mathbb{R}^n$. Suponga que $H$ tiene dos bases ortonormales, $\{ u_1,u_2, \dots, u_k \}$ y $\{ w_1,w_2, \dots, w_k \}$. Sea $v \in \mathbb{R}^n$. Entonces: $(v \cdot u_1)u_1 + (v \cdot u_2)u_2 + \cdots + (v \cdot u_k)u_k = (v \cdot w_1)w_1 + (v \cdot w_2)w_2 + \cdots + (v \cdot w_k)w_k$.
%\end{thm} 

\begin{defin}
Si V es un K espacio vectorial con producto interno $<, > $, $W$ subespacio de $V$, entonces $\|v-P_W(v)\|$, es la distancia del vector al subespacio $W$. Usaremos la notación $d(v,W)=\|v-P_W(v)\|$.
\end{defin}


Con las propiedades obtenidas para la proyección ortogonal, la distancia entre un vector y un
subespacio hereda naturalmente la siguiente propiedad:
$$d(v,W)=0 \Longleftrightarrow P_W(v)=v.$$

\begin{defin}
Sea $H$ un subespacio de $\mathbb{R}^n$. El complemento ortogonal de $H$ denotado por $H^\perp$, está dado por $H^\perp = \{ x \in \mathbb{R}^n: x \cdot h =0$ para toda $h \in H \}$.
\end{defin}


\begin{thm}
Si $H$ es un subespacio de $\mathbb{R}^n$, entonces

\begin{enumerate}
\item
$H^\perp$ es un subespacio de $\mathbb{R}^n$.

\item
$H \cap H^\perp = \{ 0 \}$.

\item
$\dim H^\perp = n-\dim H$.

\end{enumerate}
\end{thm}

\begin{defin}
Sea $H$ un subespacio de $\mathbb{R}^n$ y sea $v \in \mathbb{R}^n$. Entonces existe un par único de vectores $h$ y $p$ tales que $h \in H$, $p \in H^\perp$ y $v = h + p$. En particular $h = P_H(v)$ y $p = P_{H^\perp}(v)$ de manera que $v = h+p = P_H(v) + P_{H^\perp}(v)$.
\end{defin}








\begin{ejer}
Sea $\pi = \left\{ \begin{pmatrix} x \\ y \\ z \end{pmatrix}: 2x-y+3z=0 \right\}$. Exprese el vector $v = \begin{pmatrix} 3 \\ -2 \\ 4 \end{pmatrix}$ en términos de $h+p$.
\end{ejer} 









%\begin{itemize}
%\item
%\textbf{Teo: }Sea $H$ un subespacio de $\mathbb{R}^n$ y sea $v$ un vector en $\mathbb{R}^n$. Entonces proy$_Hv$ es la mejor aproximación para $v$ en $H$ en el siguiente sentido: si $h$ es cualquier otro vector en $H$, entonces $|v -$proy$_Hv | < |v - h|$.

%\end{itemize}

\subsection{Proceso de Ortonormalización de Gram-Schmidt: }
Dada una base de un espacio vectorial, podemos obtener una base ortonormal. A este proceso se le conoce como Ortonormalización de Gram-Schmidt.

Sea $W$ un subespacio de dimensión $m$ de $V$. Entonces $H$ tiene una base ortonormal.

\begin{itemize}

\item
\textbf{Procedimiento para obtener una base ortonormal}
Sea $B = \{v_i\}_{i=1}^m $ una base ordenada cualquiera de $W $.
\begin{enumerate}
\item Sea $w_1 = v_1$. Construyamos $w_2 = v_2 - c_1v_1$, con $c_1 =\dfrac{<v_2,v_1>}{<v_1,v_1>}, $ coeficiente de Fourier de $v_2$ sobre $v_1$, de modo que $w_2 \perp v_1$.
\item Suponemos que esta construcción de $w_k$ es posible hacerla hasta
obtener el conjunto ortogonal $ \{w_i\}_{i=1}^{m-1}$ con $w_k=v_k-\sum_{i=1}^{k-1}c_iw_i$.
\item Obtenemos $w_m=v_m-\sum_{i=1}^{m-1}c_iw_i$.
De este modo, queda demostrado que hemos construido una base ortogonal
para $W$.

\end{enumerate}
Si al final del proceso se normaliza cada vector obteniendo los nuevos
vectores $w_k$, se construirá una base ortonormal para $W$.
\end{itemize}








\begin{ejem}Construya una base ortonormal comenzando con la base 
$$v_1 = \begin{pmatrix} 1 \\ -1 \\ -1 \\ 1 \end{pmatrix}, v_2 = \begin{pmatrix} 2 \\ 1 \\ 0 \\ 1 \end{pmatrix}, v_3 = \begin{pmatrix} 2 \\ 2 \\ 1 \\ 2 \end{pmatrix}  $$
\end{ejem}

\begin{ejem} Encuentre una base ortonormal para el conjunto de vectores en $\mathbb{R}^3$ que está sobre el plano 
$$\pi = \left\{ \begin{pmatrix} x \\ y \\ z \end{pmatrix}: 2x-y+3z=0 \right\}$$
\end{ejem}